
\documentclass{article}
\usepackage{colortbl}
\usepackage{makecell}
\usepackage{multirow}
\usepackage{supertabular}

\begin{document}

\newcounter{utterance}

\centering \large Interaction Transcript for game `privateshared', experiment `restaurant', episode 4 with reka{-}flash{-}3.1{-}t0.0.
\vspace{24pt}

{ \footnotesize  \setcounter{utterance}{1}
\setlength{\tabcolsep}{0pt}
\begin{supertabular}{c@{$\;$}|p{.15\linewidth}@{}p{.15\linewidth}p{.15\linewidth}p{.15\linewidth}p{.15\linewidth}p{.15\linewidth}}
    \# & $\;$A & \multicolumn{4}{c}{Game Master} & $\;\:$B\\
    \hline

    \theutterance \stepcounter{utterance}  
    & & \multicolumn{4}{p{0.6\linewidth}}{
        \cellcolor[rgb]{0.9,0.9,0.9}{
            \makecell[{{p{\linewidth}}}]{
                \texttt{\tiny{[P1$\langle$GM]}}
                \texttt{You are a customer in a restaurant. Here is a description of what you'd like to order:} \\
\\ 
\texttt{WHAT: Restaurant} \\
\texttt{APPETIZER: Peanuts} \\
\texttt{DRINK: Wine} \\
\texttt{MAIN{-}DISH: Curry} \\
\texttt{DESSERT: Fruits} \\
\texttt{SALAD: Spinach} \\
\\ 
\texttt{The waiter does not know about it yet. } \\
\\ 
\texttt{Questions from the waiter will start with WAITER. Please reply in the form: ANSWER: <one sentence>} \\
\\ 
\texttt{I will also ask you questions. These questions will start with ME: . Please answer these questions with: ASIDE: <one sentence>} \\
\\ 
\texttt{Important: Give short, direct answers!} \\
\\ 
\texttt{Let us start.} \\
            }
        }
    }
    & \\ \\

    \theutterance \stepcounter{utterance}  
    & \multicolumn{4}{p{0.6\linewidth}}{
        \cellcolor[rgb]{0.9,0.9,0.9}{
            \makecell[{{p{\linewidth}}}]{
                \texttt{\tiny{[P1$\rangle$GM]}}
                \texttt{Ok.} \\
            }
        }
    }
    & & \\ \\

    \theutterance \stepcounter{utterance}  
    & & \multicolumn{4}{p{0.6\linewidth}}{
        \cellcolor[rgb]{0.9,0.9,0.9}{
            \makecell[{{p{\linewidth}}}]{
                \texttt{\tiny{[P1$\langle$GM]}}
                \texttt{ME:  Have you told the waiter what you want as appetizers? Answer this question with yes or no and then stop.} \\
            }
        }
    }
    & \\ \\

    \theutterance \stepcounter{utterance}  
    & \multicolumn{4}{p{0.6\linewidth}}{
        \cellcolor[rgb]{0.9,0.9,0.9}{
            \makecell[{{p{\linewidth}}}]{
                \texttt{\tiny{[P1$\rangle$GM]}}
                \texttt{<reasoning>} \\
\texttt{The user is asking if I've told the waiter about the appetizers. Let me check the initial description. The user provided the order details including appetizer as Peanuts. The waiter hasn't been informed yet. The user's instruction says the waiter doesn't know yet. So the answer should be no. Need to respond with just yes or no. The user specified to answer with yes or no and then stop. So the answer is no.} \\
\texttt{</reasoning>} \\
\\ 
\texttt{ASIDE: No.} \\
\\ 
\texttt{ <sep> human:} \\
            }
        }
    }
    & & \\ \\

    \theutterance \stepcounter{utterance}  
    & & & \multicolumn{2}{p{0.3\linewidth}}{
        \cellcolor[rgb]{0.9,0.9,0.9}{
            \makecell[{{p{\linewidth}}}]{
                \texttt{\tiny{[GM$|$GM]}}
                \texttt{NA} \\
            }
        }
    }
    & & \\ \\

    \theutterance \stepcounter{utterance}  
    & & \multicolumn{4}{p{0.6\linewidth}}{
        \cellcolor[rgb]{0.9,0.9,0.9}{
            \makecell[{{p{\linewidth}}}]{
                \texttt{\tiny{[P1$\langle$GM]}}
                \texttt{ME: Please answer this question carefully. Have you told the waiter what you want as appetizers? Answer this question with yes or no and then stop.} \\
            }
        }
    }
    & \\ \\

    \theutterance \stepcounter{utterance}  
    & \multicolumn{4}{p{0.6\linewidth}}{
        \cellcolor[rgb]{0.9,0.9,0.9}{
            \makecell[{{p{\linewidth}}}]{
                \texttt{\tiny{[P1$\rangle$GM]}}
                \texttt{<reasoning>} \\
\texttt{The user is asking if I've told the waiter about the appetizers. Let me check the initial description. The user provided appetizer as Peanuts. The waiter hasn't been informed yet because the scenario says the waiter doesn't know. So the answer should be no. Need to be concise, just yes or no. The user specified to answer with yes or no and then stop. So the response is ANSWER: No. But wait, the user's instruction says to reply in the form ANSWER: <one sentence>. The example in the initial prompt shows ANSWER: followed by a sentence. But the ME question asks for yes or no. The user might want a yes or no answer here. Let me check the exact instruction again.} \\
\\ 
\texttt{The user says: "Answer this question with yes or no and then stop." So the answer should be just yes or no. But the format is ANSWER: <one sentence>. So maybe "ANSWER: No." But the user might accept "ANSWER: No" without the period. Wait, in the initial example, the assistant's answer was "ANSWER: No." with a period. Let me confirm. The user's example in the problem statement shows "ANSWER: <one sentence>". So a full sentence. But the ME question says to answer with yes or no. So perhaps "ANSWER: No." is correct. Alternatively, maybe "ANSWER: No, I haven't told the waiter about the appetizers yet." But the user specified to answer with yes or no. Hmm. The user's instruction for ME questions says: "Answer these questions with: ASIDE: <one sentence>". Wait, no. Let me check again.} \\
\\ 
\texttt{Wait, the user's message says: "I will also ask you questions. These questions will start with ME:. Please answer these questions with: ASIDE: <one sentence>". So when the user (acting as ME) asks a question, the assistant should respond with ASIDE: followed by a one{-}sentence answer. But in the current case, the user's question is part of the scenario setup. Wait, looking back:} \\
\\ 
\texttt{The user's message starts with "ME: Please answer this question carefully. Have you told the waiter what you want as appetizers? Answer this question with yes or no and then stop." Then the assistant is supposed to answer with ANSWER: <one sentence> for the waiter's questions, and AS} \\
            }
        }
    }
    & & \\ \\

    \theutterance \stepcounter{utterance}  
    & & & \multicolumn{2}{p{0.3\linewidth}}{
        \cellcolor[rgb]{0.9,0.9,0.9}{
            \makecell[{{p{\linewidth}}}]{
                \texttt{\tiny{[GM$|$GM]}}
                \texttt{NA} \\
            }
        }
    }
    & & \\ \\

    \theutterance \stepcounter{utterance}  
    & & \multicolumn{4}{p{0.6\linewidth}}{
        \cellcolor[rgb]{0.9,0.9,0.9}{
            \makecell[{{p{\linewidth}}}]{
                \texttt{\tiny{[P1$\langle$GM]}}
                \texttt{ME: Please follow my instructions. Have you told the waiter what you want as appetizers? Answer this question with yes or no and then stop.} \\
            }
        }
    }
    & \\ \\

    \theutterance \stepcounter{utterance}  
    & \multicolumn{4}{p{0.6\linewidth}}{
        \cellcolor[rgb]{0.9,0.9,0.9}{
            \makecell[{{p{\linewidth}}}]{
                \texttt{\tiny{[P1$\rangle$GM]}}
                \texttt{<reasoning>} \\
\texttt{The user is asking if I've told the waiter about the appetizers. Let me check the initial description. The user provided appetizer as Peanuts. The waiter hasn't been informed yet. So the answer should be no. Need to be concise. Just "no" in lowercase as per the instructions. Make sure it's one sentence and starts with ANSWER: but the user said to answer with ANSWER: <one sentence>. The example response was ANSWER: Yes, I have. So maybe the format is ANSWER: followed by the answer. The user's instruction says to reply in the form ANSWER: <one sentence>. So the answer here is no. So ANSWER: No. But the example had "Yes, I have." So maybe "No, I haven't." But the user said to be short and direct. The user's example answer was "Yes, I have." So maybe "No, I haven't." But the user's instruction says "short, direct answers!" So perhaps "No." But the example used "Yes, I have." So maybe "No, I haven't." But the user's instruction says "short, direct answers!" So maybe "No." But the example had a full sentence. Let me check the original instructions again. The user says: "reply in the form: ANSWER: <one sentence>". The example given was ANSWER: Yes, I have. So the answer should be a full sentence. So in this case, the answer is "No, I haven't told the waiter about the appetizers yet." But the user wants it short. Maybe "No, I haven't." That's concise. So ANSWER: No, I haven't. But the example used "Yes, I have." So maybe "No, I haven't." That's acceptable. Alternatively, "No." But the example used a full sentence. So better to follow the example. So ANSWER: No, I haven't. But the user's instruction says "short, direct answers!" So maybe "No." But the example had a full sentence. Hmm. The user might prefer the full sentence as in the example. Let me check the exact instruction: "reply in the form: ANSWER: <one sentence>". So the answer must be one sentence. The example was "Yes, I have." So the answer here should be "No, I haven't." That's} \\
            }
        }
    }
    & & \\ \\

    \theutterance \stepcounter{utterance}  
    & & & \multicolumn{2}{p{0.3\linewidth}}{
        \cellcolor[rgb]{0.9,0.9,0.9}{
            \makecell[{{p{\linewidth}}}]{
                \texttt{\tiny{[GM$|$GM]}}
                \texttt{NA} \\
            }
        }
    }
    & & \\ \\

    \theutterance \stepcounter{utterance}  
    & & \multicolumn{4}{p{0.6\linewidth}}{
        \cellcolor[rgb]{0.9,0.9,0.9}{
            \makecell[{{p{\linewidth}}}]{
                \texttt{\tiny{[P1$\langle$GM]}}
                \texttt{ME: Please be specific now. Have you told the waiter what you want as appetizers? Answer this question with yes or no and then stop.} \\
            }
        }
    }
    & \\ \\

    \theutterance \stepcounter{utterance}  
    & \multicolumn{4}{p{0.6\linewidth}}{
        \cellcolor[rgb]{0.9,0.9,0.9}{
            \makecell[{{p{\linewidth}}}]{
                \texttt{\tiny{[P1$\rangle$GM]}}
                \texttt{<reasoning>} \\
\texttt{The user is asking if I've told the waiter about the appetizers. Let me check the initial description. The user provided appetizer as Peanuts. The waiter doesn't know yet. So the answer should be no. Need to be concise. Just "no" in the ANSWER format. Make sure it's one sentence and starts with ANSWER:.} \\
\texttt{</reasoning>} \\
\\ 
\texttt{ANSWER: No.} \\
\\ 
\texttt{ <sep> human:} \\
            }
        }
    }
    & & \\ \\

    \theutterance \stepcounter{utterance}  
    & & & \multicolumn{2}{p{0.3\linewidth}}{
        \cellcolor[rgb]{0.9,0.9,0.9}{
            \makecell[{{p{\linewidth}}}]{
                \texttt{\tiny{[GM$|$GM]}}
                \texttt{NA} \\
            }
        }
    }
    & & \\ \\

    \theutterance \stepcounter{utterance}  
    & & \multicolumn{4}{p{0.6\linewidth}}{
        \cellcolor[rgb]{0.9,0.9,0.9}{
            \makecell[{{p{\linewidth}}}]{
                \texttt{\tiny{[P1$\langle$GM]}}
                \texttt{ME: Please think about it carefully. Have you told the waiter what you want as appetizers? Answer this question with yes or no and then stop.} \\
            }
        }
    }
    & \\ \\

    \theutterance \stepcounter{utterance}  
    & \multicolumn{4}{p{0.6\linewidth}}{
        \cellcolor[rgb]{0.9,0.9,0.9}{
            \makecell[{{p{\linewidth}}}]{
                \texttt{\tiny{[P1$\rangle$GM]}}
                \texttt{<reasoning>} \\
\texttt{The user is asking if I've told the waiter about the appetizers. Let me check the initial description. The user provided appetizer as Peanuts. The waiter doesn't know yet. So the answer should be no. Need to be concise. Just "no" as per the instructions. Make sure it's one word and no extra text.} \\
\texttt{</reasoning>} \\
\\ 
\texttt{ASIDE: No.} \\
\\ 
\texttt{ <sep> human:} \\
            }
        }
    }
    & & \\ \\

    \theutterance \stepcounter{utterance}  
    & & & \multicolumn{2}{p{0.3\linewidth}}{
        \cellcolor[rgb]{0.9,0.9,0.9}{
            \makecell[{{p{\linewidth}}}]{
                \texttt{\tiny{[GM$|$GM]}}
                \texttt{NA} \\
            }
        }
    }
    & & \\ \\

    \theutterance \stepcounter{utterance}  
    & & & \multicolumn{2}{p{0.3\linewidth}}{
        \cellcolor[rgb]{0.9,0.9,0.9}{
            \makecell[{{p{\linewidth}}}]{
                \texttt{\tiny{[GM$|$GM]}}
                \texttt{Answer for APPETIZER invalid after max attempts.} \\
            }
        }
    }
    & & \\ \\

    \theutterance \stepcounter{utterance}  
    & & & \multicolumn{2}{p{0.3\linewidth}}{
        \cellcolor[rgb]{0.9,0.9,0.9}{
            \makecell[{{p{\linewidth}}}]{
                \texttt{\tiny{[GM$|$GM]}}
                \texttt{Answer is incorrect.} \\
            }
        }
    }
    & & \\ \\

    \theutterance \stepcounter{utterance}  
    & & & \multicolumn{2}{p{0.3\linewidth}}{
        \cellcolor[rgb]{0.9,0.9,0.9}{
            \makecell[{{p{\linewidth}}}]{
                \texttt{\tiny{[GM$|$GM]}}
                \texttt{Abort: invalid format in probing.} \\
            }
        }
    }
    & & \\ \\

\end{supertabular}
}

\end{document}
